% Test of hypertex html: specials, in particular when the opening special
% is given in vertical mode.  The link rectangle must cover the link text,
% and a link broken across lines must get one rectangle per line.
% Public domain.
\nopagenumbers
\parskip=1pc
\font\rm=cmr10 \rm

% --- single-line anchors ---

% opened in vertical mode, at the very start of the page
\special{html:<a href="https://tug.org/">}vmode top\special{html:</a>}

Click:\par
% opened in vertical mode, after a \par
\special{html:<a href="https://tug.org/">}vmode\special{html:</a>}

% opened in horizontal mode; this case was always correct
Click: \special{html:<a href="https://tug.org/">}hmode\special{html:</a>}

\vbox{% opened in vertical mode, at the start of a \vbox
\special{html:<a href="https://tug.org/">}vbox\special{html:</a>}
}

\hbox{% opened in horizontal mode, at the start of an \hbox
\special{html:<a href="https://tug.org/">}hbox\special{html:</a>}
}

% --- anchors whose text is broken across lines ---

\hsize=2.5in

hmode start, wraps.
\special{html:<a href="https://tug.org/">}anchor text long enough to be
broken by TeX across two lines\special{html:</a>} done.

vmode start, wraps.\par
\special{html:<a href="https://tug.org/">}anchor text long enough to be
broken by TeX across two lines\special{html:</a>} done.

vmode start in a vbox, wraps.\par
\noindent\vbox{\hsize=2.5in
\special{html:<a href="https://tug.org/">}anchor text long enough to be
broken by TeX across two lines\special{html:</a>} done.}

\bye
